% =============================================================================
%  lemma_88_short_body.tex -- a short proof of the extension lemma, after
%  Deligne's tame case (following Tendler and Guignard).
%
%  Self-contained: labels defined here carry the suffix "_short"; all other
%  \Cref's resolve against ../main.aux. If this proof is adopted, rename
%  lemma:ExtensionOfTamelyRamifiedSheafShort back to
%  lemma:ExtensionOfTamelyRamifiedSheaf and \input this file in place of the
%  placeholder in sections/GeometricCFT/ProofOfReducedTheorem.tex.
% =============================================================================

\subsection*{The extension lemma, after Deligne's tame case}

We keep the notation of \Cref{section:BlowupOfSymmetricPowerOfCurves}:
$C$ is a projective smooth geometrically connected curve over the perfect field $k$,
$\mdls = \sum n_i P_i > 0$ is a modulus with complement $U = C \setminus \mdls$,
and $d$ is an integer satisfying \Cref{equation:degree-condition}.
We write $Z_0 = Z_0(\mdls, d) \subset C^{(d)}$ for the center of the blowup
$\pi : X_{\mdls, d} \to C^{(d)}$, $E_0$ for its exceptional divisor with generic point $\eta_0$,
$\Cmod{d} \subseteq X_{\mdls,d}$ for Takeuchi's compactification
(\Cref{theorem:takeuchi-compactification-theorem}), and
$H = \Cmod{d} \setminus U^{(d)} = E_0 \times_{C^{(d)}} \Cmod{d}$ for the boundary divisor.
Throughout, $p$ denotes the characteristic of $k$; if $p = 0$, statements of the form
``prime to $p$'' are to be read as no condition, and the arguments below only simplify.
We set
\[
N \;:=\; d - \deg\mdls - g + 1 \;\geq\; 1,
\]
the relative dimension of $\Phi_d$ (\Cref{section:AbelJacobi}; positivity
follows from \Cref{equation:degree-condition}).

\begin{lemma}\label{lemma:ExtensionOfTamelyRamifiedSheafShort}
    If $\cF^{[d]}$ is a locally constant sheaf on $U^{(d)}$ which is tamely ramified along the
    boundary divisor $H = \tilde{C}^{(d)}_\mdls \setminus U^{(d)}$,
    then $\cF^{[d]}$ extends to a locally constant sheaf $\tilde{\cF}^{[d]}$ on
    $\tilde{C}^{(d)}_\mdls$.
\end{lemma}

The proof has the same two-step skeleton as the other routes --- by
Zariski--Nagata purity it suffices to show the corresponding torsor is
\emph{unramified} at $\eta_0$ (\Cref{prop:extension_via_purity_short}), and
tameness at $\eta_0$ is then upgraded to unramifiedness
(\Cref{prop:tame_implies_unramified_short}) --- but the second step is now
soft. The point $\eta_0$ lies in the \emph{generic fiber} of the compactified
Abel--Jacobi morphism $\tilde{\Phi}_d$, with the same local ring; the
corresponding generic fiber of the \emph{open} morphism $\Phi_d$ is an honest
affine space $\A^N_\kappa$ over the function field $\kappa$ of $\PicCm[d]$
(\Cref{lemma:boundary_generic_fiber_short}). Since affine space has no
nontrivial abelian covers of degree prime to $p$
(\Cref{lemma:affine_pro_p_short} --- five lines of Kummer theory), the
prime-to-$p$ part of our torsor is \emph{constant} on that fiber, hence
unramified at $\eta_0$; and the $p$-primary part of the inertia is killed by
tameness alone, since a Galois extension of $p$-power order in residue
characteristic $p$ admits no nontrivial tame ramification. No valuation
computation and no degree count is needed.

\subsubsection*{Geometry of the boundary and of the generic fiber}

Let $\xi$ denote the generic point of $\PicCm[d]$ and let
$\kappa := \kappa(\xi)$ be its residue field, i.e.\ the function field of
$\PicCm[d]$.

\begin{lemma}\label{lemma:boundary_generic_fiber_short}
    With notation as above:
    \begin{enumerate}
        \item $\Cmod{d}$ is an integral scheme, smooth over $k$; $H$ is a nonempty open
              subscheme of $E_0$, irreducible of codimension $1$ in $\Cmod{d}$ with generic
              point $\eta_0$; the point $\eta_0$ is the unique point of
              $\Cmod{d} \setminus U^{(d)}$ whose local ring has dimension $1$, and
              $\Oo := \Oo_{\Cmod{d}, \eta_0}$ is a discrete valuation ring with fraction
              field $K := k(C^{(d)})$.
        \item $\tilde{\Phi}_d(\eta_0) = \xi$. The generic fiber
              $P_\xi := \tilde{\Phi}_d^{-1}(\xi)$ is isomorphic to $\bP^N_\kappa$, its
              function field is $K$, and the local ring of $P_\xi$ at $\eta_0$ \emph{equals}
              $\Oo$. In particular $\kappa \subseteq \Oo$.
        \item $A_\xi := P_\xi \cap U^{(d)} = \Phi_d^{-1}(\xi)$ is a dense open subscheme of
              $P_\xi$ with closed complement $P_\xi \cap H$, its generic point is
              $\Spec(K)$, and
              \[ A_\xi \;\cong\; \A^N_\kappa \]
              as a $\kappa$-scheme.
    \end{enumerate}
\end{lemma}
\begin{proof}
    (1) $\Cmod{d}$ is an open subscheme of the blowup $X_{\mdls,d}$ of the integral scheme
    $C^{(d)}$, hence integral, with function field $K$; it is a projective space bundle over
    the smooth $k$-scheme $\PicCm[d]$ (\Cref{theorem:takeuchi-compactification-theorem}),
    hence smooth over $k$. The base change
    $H = E_0 \times_{X_{\mdls,d}} \Cmod{d} \hookrightarrow E_0$ of the open immersion
    $\Cmod{d} \hookrightarrow X_{\mdls,d}$ is an open immersion, and $H \neq \emptyset$:
    otherwise $\Phi_d = \tilde{\Phi}_d$ would be proper, while its fibers are affine spaces
    of dimension $N \geq 1$ (\Cref{section:AbelJacobi}). A nonempty open subscheme of the
    irreducible $E_0$ is irreducible with the same generic point $\eta_0$; every other point
    of $H$ is a proper specialization of $\eta_0$, so has local ring of dimension $\geq 2$.
    Finally $\Oo$ is a one-dimensional regular local ring, i.e.\ a DVR, with fraction field
    $k(\Cmod{d}) = K$.

    (2) Every fiber of $\tilde{\Phi}_d$ meets $H$: a fiber contained in $U^{(d)}$ would be a
    proper scheme equal to a fiber of $\Phi_d$, which is a positive-dimensional affine space
    over the separable closure --- not proper. Hence $\tilde{\Phi}_d(H) = \PicCm[d]$
    ($\tilde{\Phi}_d$ is proper, so the image of the closed set $H$ is closed and, as just
    shown, it meets every point). Since $H = \overline{\{\eta_0\}} \cap \Cmod{d}$,
    \[
    \PicCm[d] = \tilde{\Phi}_d(H) \subseteq
    \overline{\{\tilde{\Phi}_d(\eta_0)\}},
    \]
    which forces $\tilde{\Phi}_d(\eta_0) = \xi$. Next, since $\xi$ is the generic point of
    the integral scheme $\PicCm[d]$, its local ring is the field $\kappa$, so
    \[
    P_\xi = \Cmod{d} \times_{\PicCm[d]} \Spec \Oo_{\PicCm[d], \xi}
          = \varprojlim_{V \ni \xi \text{ open}} \tilde{\Phi}_d^{-1}(V)
    \]
    is a filtered limit of \emph{open} subschemes of $\Cmod{d}$ along the identity maps;
    consequently, for every point $y \in P_\xi$ the local rings agree,
    $\Oo_{P_\xi, y} = \Oo_{\Cmod{d}, y}$. Applying this to $y = \eta_0$ gives
    $\Oo_{P_\xi,\eta_0} = \Oo$, and in particular $\kappa = \Oo_{\PicCm[d],\xi} \subseteq \Oo$.
    By \Cref{theorem:takeuchi-compactification-theorem}, $\Cmod{d} = \bP(\cE)$ for a locally
    free sheaf $\cE$ on $\PicCm[d]$, so $P_\xi = \bP(\cE_\xi) \cong \bP^{N}_\kappa$ (the
    fiber dimension is $N$ because $P_\xi$ contains the fiber of $\Phi_d$, of dimension $N$,
    as a dense open subset --- see (3)). The generic point of $\Cmod{d}$ maps to $\xi$
    (dominance) and lies in $P_\xi$ as its generic point, so $\kappa(P_\xi) = K$.

    (3) $P_\xi \cap U^{(d)} = P_\xi \setminus H$ holds because
    $\Cmod{d} \setminus U^{(d)} = H$ (\Cref{theorem:takeuchi-compactification-theorem}), and
    this set equals $\Phi_d^{-1}(\xi)$ because $\tilde{\Phi}_d$ extends $\Phi_d$. It is open
    in the irreducible $P_\xi$ and nonempty ($\Phi_d$ is surjective,
    \Cref{section:AbelJacobi}), hence dense with the same generic point $\Spec K$.

    For the last isomorphism, recall the modular description of the fibers of $\Phi_d$
    (\cite[Section 4, proof of Theorem 4.14]{Guignard2018}; it is also the description
    recalled in \Cref{section:AbelJacobi}). Since $\mdls_\kappa \to \Spec \kappa$ is finite
    flat and surjective, the rigidified Picard presheaf is already a sheaf, so the
    $\kappa$-point $\xi$ of $\PicCm[d]$ is represented by an actual pair $(\cL, \psi)$ on
    $C_\kappa$. For a $\kappa$-scheme $T$, a $T$-point of $\Phi_d^{-1}(\xi)$ is a divisor
    $D \in U^{(d)}(T)$ together with the (unique, since rigidified pairs have no nontrivial
    automorphisms) isomorphism $(\Oo_C(D), \psi_D) \cong (\cL, \psi)_T$; sending $D$ to the
    image $s \in \Gamma(C_T, \cL_T)$ of the canonical section $1_D$ identifies
    $\Phi_d^{-1}(\xi)(T)$ with
    \[
    V_\psi(T) := \{\, s \in \Gamma(C_T, \cL_T) \;:\; s|_{\mdls_T} = \psi_T(1) \,\}
    \]
    (the condition $s|_{\mdls_T} = \psi_T(1)$ pins down the scaling, and such an $s$ is
    fiberwise nonzero because it is invertible along $\mdls$, so $D = V(s)$ is a relative
    effective Cartier divisor of degree $d$ prime to $\mdls$). By the degree condition, on
    each fiber $H^1(C, \cL(-\mdls)) = 0$, so (cohomology and base change, as recalled in
    \Cref{section:BlowupOfSymmetricPowerOfCurves}) the restriction map of vector groups
    $\mathrm{pr}_*\cL \to \mathrm{pr}_*(\cL|_{\mdls})$ over $\kappa$ is surjective with
    kernel the vector group of $W := H^0(C_\kappa, \cL(-\mdls))$, of dimension $N$
    (Riemann--Roch). Hence $A_\xi \cong V_\psi$ is a torsor under the vector group
    $\mathbb{W} \cong \bG_a^N$ over $\kappa$. Every such torsor is trivial, since
    $H^1_{et}(\Spec \kappa, \bG_a^N) = H^1(\mathrm{Gal}(\kappa^{sep}/\kappa),
    (\kappa^{sep})^N) = 0$ (additive Hilbert 90). Therefore
    $A_\xi \cong \mathbb{W} \cong \A^N_\kappa$.
\end{proof}

\subsubsection*{Affine space has no prime-to-$p$ abelian covers}

The following is the elementary engine of Deligne's tame case, as used in
(\cite{tendler2015geometricclassfieldtheory}, last section:
``$\pi_1(\A^n)^{ab}$ is a pro-$p$ group'') and in
(\cite{Guignard2018}, Section 5: triviality of the tame fundamental group of
the affine line).

\begin{lemma}\label{lemma:affine_pro_p_short}
    Let $F$ be a separably closed field and let $G'$ be a finite abelian group whose order
    is invertible in $F$. Then every $G'$-torsor on $(\A^N_F)_{et}$ is trivial.
    Equivalently, every quotient of $\pi_1^{et,ab}(\A^N_F)$ of order invertible in $F$
    vanishes (for $\operatorname{char} F = p > 0$: $\pi_1^{et,ab}(\A^N_F)$ is pro-$p$).
\end{lemma}
\begin{proof}
    Writing $G' \cong \prod_i \Z/n_i\Z$ and decomposing the torsor accordingly
    (\Cref{prop:quotient_torsors}), we may assume $G' = \Z/n\Z$ with $n$ invertible in $F$.
    Since $F$ is separably closed and $x^n - 1$ is separable, $F$ contains all $n$-th roots
    of unity, so $\Z/n\Z \cong \mu_n$ as étale sheaves on $\A^N_F$, and it suffices to show
    $H^1_{et}(\A^N_F, \mu_n) = 0$ (étale $H^1$ with coefficients in a finite abelian
    constant group classifies torsors under it, cf.\
    \Cref{cor:abelian_equivilance_fundamental_torsors}). The Kummer sequence
    $1 \to \mu_n \to \bG_m \xrightarrow{n} \bG_m \to 1$, exact on the étale site because
    $n$ is invertible, gives
    \[
    \Gamma(\A^N_F, \Oo)^\times \xrightarrow{\;n\;} \Gamma(\A^N_F, \Oo)^\times
    \longrightarrow H^1_{et}(\A^N_F, \mu_n) \longrightarrow \Pic(\A^N_F)[n].
    \]
    Now $\Gamma(\A^N_F, \Oo)^\times = F[x_1, \dots, x_N]^\times = F^\times$, which is
    $n$-divisible ($x^n - a$ is separable for $a \neq 0$, and $F$ is separably closed), and
    $\Pic(\A^N_F) = 0$ because $F[x_1,\dots,x_N]$ is a unique factorization domain. Hence
    $H^1_{et}(\A^N_F, \mu_n) = 0$.
\end{proof}

\subsubsection*{Step 1: purity, and extending unramified torsors}

This step is identical to the one in the other two routes; we reproduce it for
completeness.

\begin{prop}\label{prop:extension_via_purity_short}
    Let $G$ be a finite abelian group, let $\rho : \pi_1^{et}(U^{(d)}, \bar{u}) \to G$ be a
    continuous homomorphism and let $\cQ$ be the corresponding $G$-torsor on
    $U^{(d)}_{\text{\'et}}$ (\Cref{cor:abelian_equivilance_fundamental_torsors}).
    If $\cQ$ is unramified over $\Cmod{d}$ at $(H, \eta_0)$ in the sense of
    \Cref{definition:tamely_ramified_in_codim_1}, then $\rho$ factors, necessarily uniquely,
    through the canonical surjection
    $\iota_* : \pi_1^{et}(U^{(d)}, \bar{u}) \twoheadrightarrow \pi_1^{et}(\Cmod{d}, \bar{u})$,
    where $\iota : U^{(d)} \hookrightarrow \Cmod{d}$ is the open immersion.
    Consequently, $\cQ$ extends to a $G$-torsor on $(\Cmod{d})_{\text{\'et}}$, uniquely up
    to isomorphism.
\end{prop}
\begin{proof}
    \textbf{Surjectivity of $\iota_*$.} For every connected finite étale cover
    $T \to \Cmod{d}$, the scheme $T$ is regular (étale over the regular $\Cmod{d}$), hence
    irreducible, and $T \times_{\Cmod{d}} U^{(d)}$ is a nonempty open subscheme of $T$,
    hence connected. If the image of $\iota_*$ were contained in a proper open subgroup
    $M \leq \pi_1^{et}(\Cmod{d}, \bar{u})$, the connected cover corresponding to $M$ would
    pull back to a disconnected cover of $U^{(d)}$ --- a contradiction.

    \textbf{Normalization.} Let $G' = \rho(\pi_1^{et}(U^{(d)}, \bar{u}))$ and let
    $V \to U^{(d)}$ be the connected Galois cover with group $G'$ corresponding to
    $\ker \rho$ (\Cref{theorem:equivalence_fundamental_covers}); as a $U^{(d)}$-scheme,
    $\cQ$ is a disjoint union of $[G:G']$ copies of $V$, so $V$ is unramified over
    $\Cmod{d}$ at $(H, \eta_0)$ as well. $V$ is integral and $K(V)/K$ is finite separable.
    Let $\nu : W \to \Cmod{d}$ be the normalization of $\Cmod{d}$ in $K(V)$; since
    $\Cmod{d}$ is Nagata (\stackstag{0335}), $\nu$ is finite, $W$ is integral and normal,
    and $\nu^{-1}(U^{(d)}) \cong V$ by uniqueness of normalization.

    \textbf{$\nu$ is étale, by purity of the branch locus.} We apply \stackstag{0BMB} at
    each point $w \in W$ over $y = \nu(w)$: $\Oo_{W,w}$ is normal, $\Oo_{\Cmod{d},y}$ is
    regular, $\nu$ is finite hence quasi-finite, and
    $\dim(\Oo_{W,w}) = \dim(\Oo_{\Cmod{d},y})$ by the dimension formula for the finite
    dominant morphism $\nu$ of integral schemes of finite type over $k$. For the
    codimension-one condition, let $w' \rightsquigarrow w$ with
    $\dim(\Oo_{W,w'}) = 1$ and $y' = \nu(w')$, so $\dim(\Oo_{\Cmod{d},y'}) = 1$. If
    $y' \in U^{(d)}$, then $\nu$ is étale at $w'$ because $\nu^{-1}(U^{(d)}) = V$ is étale
    over $U^{(d)}$. Otherwise $y' = \eta_0$ by
    \Cref{lemma:boundary_generic_fiber_short}(1), and $\Oo_{W,w'}$ is a localization of
    the integral closure of the discrete valuation ring $\Oo$ in $K(V)$; the hypothesis
    that $V$ is unramified at $(H, \eta_0)$ says precisely
    (\Cref{definition:tamely_ramified_dvrs}) that the ramification index there is $1$ with
    separable residue extension, i.e.\ that $\nu$ is unramified at $w'$. Thus \stackstag{0BMB}
    applies and $\nu$ is finite étale.

    \textbf{Factorization.} As $V = W \times_{\Cmod{d}} U^{(d)}$, the action of
    $\pi_1^{et}(U^{(d)}, \bar{u})$ on the fiber $V_{\bar{u}} = W_{\bar{u}}$ is the
    restriction along $\iota_*$ of the $\pi_1^{et}(\Cmod{d}, \bar{u})$-action; the kernel
    of the former is $\ker \rho$, so $\ker(\iota_*) \subseteq \ker \rho$ and $\rho$ factors
    as $\rho = \tilde{\rho} \circ \iota_*$, uniquely by surjectivity of $\iota_*$. The
    $G$-torsor $\tilde{\cQ}$ corresponding to $\tilde{\rho}$ restricts to $\cQ$ on
    $U^{(d)}$, and any two extensions have equal characters, hence are isomorphic.
\end{proof}

\subsubsection*{Step 2: at the boundary, tame implies unramified}

\begin{prop}\label{prop:tame_implies_unramified_short}
    Let $G$ be a finite abelian group and let $\cQ$ be a $G$-torsor on
    $U^{(d)}_{\text{\'et}}$. If $\cQ$ is tamely ramified over $\Cmod{d}$ at $(H, \eta_0)$,
    then $\cQ$ is unramified at $(H, \eta_0)$.
\end{prop}

Note that no bound on the ramification of $\cQ$ along $\mdls$ is assumed, and that
tameness is used only to kill the $p$-primary part of the inertia: the prime-to-$p$ part
dies unconditionally, by the geometry of the generic fiber.

\begin{proof}
    Let $\widehat{K} = \widehat{K}_{\eta_0}$ be the fraction field of the completion of the
    discrete valuation ring $\Oo$ (\Cref{lemma:boundary_generic_fiber_short}(1)), and let
    \[ \rho : G_{\widehat{K}} \longrightarrow G \]
    be the continuous homomorphism classifying the restriction of $\cQ$ to
    $\Spec(\widehat{K})$. Let $\widehat{K}^{ur} \subseteq \widehat{K}^{sep}$ be the maximal
    unramified extension and $I := \mathrm{Gal}(\widehat{K}^{sep}/\widehat{K}^{ur})$ the
    inertia subgroup; the assertion to be proved is $\rho(I) = \{1\}$
    (\Cref{definition:tamely_ramified_dvrs}).

    \textbf{Tameness confines the inertia to the prime-to-$p$ part.}
    Write $G = G_p \times G'$ with $G_p$ a $p$-group and $|G'|$ prime to $p$. The fixed
    field of $\ker \rho$ is a finite Galois extension $F/\widehat{K}$ with group
    $\rho(G_{\widehat{K}}) \subseteq G$ and inertia subgroup $\rho(I)$. By hypothesis this
    extension is tame: its residue extension is separable and its ramification index $e$
    is prime to $p$. For such extensions $|\rho(I)| = e$ (\cite[Ch.\ IV]{serreLF}), so
    $\rho(I)$ is a finite group of order prime to $p$, whence
    $\rho(I) \subseteq \{1\} \times G'$. Letting $\rho' := \mathrm{pr}_{G'} \circ \rho$,
    which classifies the quotient $G'$-torsor $\cQ' := \cQ / G_p$
    (\Cref{prop:quotient_torsors}), it therefore suffices to prove $\rho'(I) = \{1\}$.

    \textbf{The prime-to-$p$ part is constant on the generic fiber.}
    By \Cref{lemma:boundary_generic_fiber_short}, $A_\xi \cong \A^N_\kappa$ is a
    subscheme of $U^{(d)}$ whose generic point is $\Spec(K)$, and
    $\kappa \subseteq \Oo \subseteq \widehat{K}$. Consider the restriction
    $\cQ'|_{A_\xi}$, a $G'$-torsor on $\A^N_\kappa$. Its pullback to
    $\A^N_{\kappa^{sep}}$ is trivial by \Cref{lemma:affine_pro_p_short} ($|G'|$ is
    invertible in $\kappa$, since it is prime to $p$). A trivializing section is a section
    of the finite étale scheme $\cQ'|_{A_\xi}$ over $\A^N_{\kappa^{sep}} =
    \varprojlim_{\kappa'} \A^N_{\kappa'}$, the limit running over the finite separable
    subextensions $\kappa'/\kappa$ of $\kappa^{sep}$; as $\cQ'|_{A_\xi}$ is of finite
    presentation, the section is defined at some finite level. Hence there is a finite
    separable extension $\kappa'/\kappa$ such that
    \[
    \cQ'|_{A_\xi \times_\kappa \kappa'} \;\text{ is the trivial } G' \text{-torsor.}
    \]

    \textbf{Constant covers are unramified.}
    Let $\widehat{\Oo}$ be the completion of $\Oo$, so $\widehat{K} =
    \Frac(\widehat{\Oo})$ and $\kappa \subseteq \Oo \subseteq \widehat{\Oo}$. Since
    $\kappa'/\kappa$ is finite separable, the base change
    $\widehat{\Oo} \otimes_\kappa \kappa'$ is finite étale over the complete discrete
    valuation ring $\widehat{\Oo}$; hence it is a finite product of complete discrete
    valuation rings, each of which is unramified over $\widehat{\Oo}$ (étale means
    exactly: the maximal ideal of $\widehat{\Oo}$ generates the maximal ideal, i.e.\ the
    ramification index is $1$, and the residue extension is separable). Inverting a
    uniformizer, $\widehat{K} \otimes_\kappa \kappa' = \prod_j \widehat{K}_j$ is a finite
    product of finite separable field extensions of $\widehat{K}$, each \emph{unramified}
    over $\widehat{K}$. Fix one factor $\widehat{K}_1$ and embed it in
    $\widehat{K}^{sep}$; then
    $\widehat{K}_1 \subseteq \widehat{K}^{ur}$, so
    \[
    I \;=\; \mathrm{Gal}(\widehat{K}^{sep}/\widehat{K}^{ur})
      \;\subseteq\; \mathrm{Gal}(\widehat{K}^{sep}/\widehat{K}_1) \;=\; G_{\widehat{K}_1}.
    \]
    Now $\Spec(\widehat{K}_1) \to \Spec(\widehat{K} \otimes_\kappa \kappa') \to
    \Spec(K \otimes_\kappa \kappa') \to A_\xi \times_\kappa \kappa'$ (the last map is the
    inclusion of a generic point, \Cref{lemma:boundary_generic_fiber_short}(3)), so the
    restriction of $\rho'$ to $G_{\widehat{K}_1}$ classifies the pullback of the
    \emph{trivial} torsor $\cQ'|_{A_\xi \times_\kappa \kappa'}$, and is therefore the
    trivial homomorphism:
    \[
    \rho'\left(G_{\widehat{K}_1}\right) = \{1\}, \qquad\text{hence}\qquad
    \rho'(I) \subseteq \rho'\left(G_{\widehat{K}_1}\right) = \{1\}.
    \]
    Combining with the first step, $\rho(I) \subseteq \{1\} \times G'$ has trivial image
    under $\mathrm{pr}_{G'}$, so $\rho(I) = \{1\}$, i.e.\ $\cQ$ is unramified at
    $(H, \eta_0)$.
\end{proof}

\subsubsection*{Conclusion of the proof}

\begin{proof}[Proof of \Cref{lemma:ExtensionOfTamelyRamifiedSheafShort}]
    Let $G = \Lambda^\times$ and let $\cP^{[d]}$ be the $G$-torsor on
    $U^{(d)}_{\text{\'et}}$ corresponding to $\cF^{[d]}$ under
    \Cref{prop:torsor_module_equivalence}, with character
    $\rho : \pi_1^{et}(U^{(d)}, \bar{u}) \to G$.
    By \Cref{lemma:boundary_generic_fiber_short}(1), the only codimension-one point of
    $\Cmod{d} \setminus U^{(d)}$ is $\eta_0$, so the hypothesis that $\cF^{[d]}$ is tamely
    ramified along $H$ means precisely that $\cP^{[d]}$ is tamely ramified over $\Cmod{d}$
    at $(H, \eta_0)$ (\Cref{definition:tamely_ramified_in_codim_1}).
    By \Cref{prop:tame_implies_unramified_short}, $\cP^{[d]}$ is unramified at
    $(H, \eta_0)$, and by \Cref{prop:extension_via_purity_short} the character $\rho$
    factors uniquely through $\pi_1^{et}(\Cmod{d}, \bar{u})$. The locally constant sheaf
    $\tilde{\cF}^{[d]}$ of free rank-one $\Lambda$-modules corresponding to the factored
    character (\Cref{cor:abelian_equivilance_fundamental_torsors} together with
    \Cref{prop:torsor_module_equivalence}) restricts to $\cF^{[d]}$ on $U^{(d)}$.
\end{proof}

\begin{rem*}
    As in the other routes, the proof also gives uniqueness of the extension, and full
    faithfulness of restriction from locally constant sheaves on $\Cmod{d}$ to locally
    constant sheaves on $U^{(d)}$, since
    $\iota_* : \pi_1^{et}(U^{(d)}) \to \pi_1^{et}(\Cmod{d})$ is surjective
    (\Cref{prop:extension_via_purity_short}). This is the source of the uniqueness
    assertions in \Cref{thm:GCFT_torsors} and \Cref{thm:GCFT_reduced}.
\end{rem*}

\begin{rem*}[Relation to Deligne's tame case, and to Routes 1 and 2]
    In Deligne's tame case (modulus $\mdls$ reduced;
    \cite{tendler2015geometricclassfieldtheory}, last section) the two soft inputs above
    are used directly on the open Abel--Jacobi fibration: a tame character becomes
    unramified on the curve after raising to a prime-to-$p$ power, the unramified case
    descends it, and pro-$p$-ness of $\pi_1^{ab}$ of the affine-space fibers transfers the
    conclusion back to the character itself. With wild ramification this reduction is
    impossible \emph{on the curve} --- no prime-to-$p$ power of a wildly ramified character
    is unramified --- which is why the thesis passes through Takeuchi's compactification
    and the tameness theorem. The proof above is the same mechanism executed once, at the
    single boundary point $\eta_0$: tameness (the thesis's hard input) kills the $p$-part
    of the inertia, and the affine-space geometry of the generic Abel--Jacobi fiber kills
    the prime-to-$p$ part. In the language of the placeholder note, this is
    \textbf{Route~1}: the kernel of
    $\pi_1^{ab}(U^{(d)}) \to \pi_1^{ab}(\PicCm[d])$ meets our character only in its
    pro-$p$ part. Route~2 (\texttt{lemma\_88\_route\_2/}) and the direct valuation
    computation (\texttt{extension\_lemma\_completion/}) prove the same statement; the
    present route replaces their Kummer-valuation and degree-count arguments by the
    triviality statement \Cref{lemma:affine_pro_p_short}.

    The global geometry is indispensable: for a proper submodulus
    $\ndls \subsetneq \mdls$ the analogous statement at $\eta_\ndls$ is \emph{false} (the
    computations of \Cref{section:ramification_after_blowup_is_tame} produce genuinely,
    tamely, ramified symmetric powers there), and correspondingly the exceptional divisor
    of the blowup at $\ndls$ does not sit inside a fibration with affine-space fibers over
    a base of the descent.
\end{rem*}
